% Source: https://artofproblemsolving.com/community/c3h2488564_root_i_of_i
% Stable identity: aops-topic-2488564-q01
\textbf{(a)} $\ln e^{\dfrac\pi2}=\dfrac\pi2$, so $x^i=\cos\dfrac\pi2+i\sin\dfrac\pi2=0+i=i$.

\textbf{(b)} We need $\cos(\ln x)=0$ and $\sin(\ln x)=1$, i.e. $\ln x=\dfrac\pi2+2k\pi$ for some integer $k$. Since $\ln$ takes every real value exactly once, the solutions are
\[\boxed{x=e^{\dfrac\pi2+2k\pi},\quad k\in\mathbb Z},\]
for example $e^{\dfrac\pi2}\approx4.810$, $e^{\dfrac{5\pi}{2}}\approx2575.97$ and $e^{-\dfrac{3\pi}{2}}\approx0.00898$. So the poster's value $e^{\dfrac\pi2}=\sqrt{e^\pi}$ works, but it is not the only one.
